%% SEGAN manuscript - CurtailmentIQ
%% Elsevier elsarticle template. Compile with pdfLaTeX in Overleaf.
%% VERSION 6.0 - reencuadre de la revision: se corrige la fecha del cambio de
%% regimen, se elimina la atribucion causal al almacenamiento, se agrega la
%% descomposicion estacional. Ningun resultado, frontera temporal ni tabla de
%% resultados cambia respecto de la version enviada (tag submitted-segan-v1).
%% Requiere las 5 figuras PDF en la carpeta del proyecto: fig1_plants_map.pdf,
%% fig2_distribution.pdf, fig3_rolling_coverage.pdf, fig4_coverage_width.pdf,
%% fig5_regime_changepoints.pdf
\documentclass[preprint,3p,times]{elsarticle}

\usepackage{amsmath,amssymb,amsfonts}
\usepackage{graphicx}
\usepackage{booktabs}
\usepackage{url}
\usepackage[hidelinks]{hyperref}
\usepackage{siunitx}
\usepackage{multirow}

\journal{Sustainable Energy, Grids and Networks}

\begin{document}

\begin{frontmatter}

\title{Probabilistic forecasting of renewable curtailment under regime change: an open dataset and a conformal method with distribution-shift-aware coverage control for Chile's National Electricity System}

\author[usach]{Pablo Reyes Cerda\corref{cor1}}
\ead{pablo.reyes.ce@usach.cl}
\author[ubb]{Kerven Cea Morales}
\ead{kceam@ubiobio.cl}

\cortext[cor1]{Corresponding author.}

\address[usach]{Mag\'ister en Gesti\'on de la Innovaci\'on y el Emprendimiento Tecnol\'ogico, Universidad de Santiago de Chile, Santiago, Chile}
\address[ubb]{Programa de Doctorado en Matem\'atica Aplicada, Universidad del B\'io-B\'io, Concepci\'on, Chile}

\begin{abstract}
The curtailment of variable renewable generation has become a first-order operational and economic problem in power systems with rapid solar and wind expansion. Chile's National Electricity System (SEN) curtailed roughly 19{,}000~GWh of renewable energy between 2022 and mid-2026, and over that period the statistical behaviour of curtailment changed in a way that breaks the assumptions of standard forecasting practice. The change has two distinct components: a gradual distributional shift in the magnitude of curtailment that accumulates between mid-2023 and late 2024, and a break in trend at the turn of 2024 to 2025, after which magnitudes stop growing even as the solar fleet keeps expanding. This paper makes two contributions. First, we release an open, plant-level, errata-corrected dataset of daily and hourly curtailment for the entire SEN, covering 53 continuous months and 300 plants across four technologies, whose annual totals reconcile exactly with the official system operator figures. Second, using this dataset, we develop a probabilistic seven-day-ahead curtailment forecaster and pair it with a conformal calibration layer designed to remain informative under distribution shift. We show that point forecasting offers little marginal value for this problem, since a seasonal-naive baseline is difficult to beat in absolute error, and that the operationally relevant challenge is therefore reliable uncertainty quantification rather than point accuracy. We further show that this regime change degrades naive conformal calibration not by breaking coverage but by inflating it: the static split remains over-covered through the transition and its intervals widen substantially. We also show that a substantial part of what reads as regime change is in fact the annual cycle: in the test window of maximum divergence from the calibration distribution, October to December 2024, between 61 and 76 per cent of the apparent jump in mean log-magnitude is seasonal rather than secular, and we separate the two components explicitly. On a synthetic panel that isolates covariate shift, the textbook weighted-conformal remedy fails by returning uninformative infinite intervals; on the real data, an adaptive scheme that combines a probability-integral-transform score with online recalibration and a transport of the calibration scores restores sharpness while holding coverage near the nominal level through the transition window, narrowing the transition intervals by about a third. The static split provides finite-sample marginal coverage under exchangeability, and the adaptive layer provides long-run coverage control when exchangeability fails. The dataset and code are publicly available.
\end{abstract}

\begin{keyword}
renewable curtailment \sep probabilistic forecasting \sep conformal prediction \sep distribution shift \sep change-point detection \sep power system operations
\end{keyword}

\end{frontmatter}

%==================================================================
\section{Introduction}
\label{sec:intro}

The share of variable renewable energy in electricity systems has grown faster than the transmission and demand able to absorb it, and curtailment, the deliberate reduction of available solar, wind or hydropower output, has become a routine consequence of that imbalance rather than an exceptional event, documented across systems worldwide~\cite{oshaughnessy2020}. In Chile's National Electricity System (Sistema El\'ectrico Nacional, SEN), recorded curtailment rose from about 1.5~TWh in 2022 to more than 6~TWh per year by 2024, and the cumulative figure over the 2022 to mid-2026 period reaches roughly 19~TWh. For the owners of affected plants this is generated energy that is never remunerated, and for the system it is a growing signal of the mismatch between where and when renewable energy is produced and where and when it can be used.

Anticipating curtailment is therefore of direct operational and commercial value. An asset owner who knows, days in advance, how much a given plant is likely to curtail can adjust maintenance, commercial and dispatch decisions accordingly. Yet curtailment is considerably harder to forecast than generation. Generation depends primarily on weather and installed capacity, which are well understood and increasingly predictable, whereas curtailment depends on the state of the whole system: net load, transmission congestion, the marginal cost of energy, and the behaviour of other assets. It is a downstream, system-level quantity, and this makes it both more difficult and less studied than generation forecasting.

Two features of the Chilean case make it a particularly informative setting. First, curtailment in the SEN is highly concentrated: a small number of regions in the north, where solar capacity has outpaced the north-to-south transmission corridor, account for the large majority of national curtailment. Second, and central to this paper, the statistical distribution of curtailment in the SEN changed materially over the study period. The change is gradual rather than abrupt: a formal change-point analysis of the monthly mean log-magnitude, described in Section~\ref{sec:data}, places the shift in a ramp of roughly eighteen months running from mid-2023 to late 2024, with no detected break in 2025 and none in October 2024. A forecasting model calibrated on one part of that trajectory is therefore miscalibrated on a later part, not because the model is wrong but because the system it describes has changed. Two features of the change make it a demanding test case rather than a textbook one. First, it has no single identifiable cause: the growth of solar capacity behind a saturated transmission corridor, the entry of battery energy storage, transmission reinforcements and changes in northern mining demand all act over the same months, and the data released here cannot separate them. Second, it is superimposed on a strong annual cycle, so that separating the seasonal component from the secular one is itself part of the problem: in the window where the divergence from the calibration distribution is largest, the majority of the apparent jump turns out to be seasonal. This is a concrete instance of the distribution-shift problem that pervades data-driven forecasting, and it is the methodological core of this work, precisely because a calibration layer must remain reliable when the cause of the shift is unknown and partly confounded with seasonality.

A further difficulty, often underappreciated in applied forecasting, is that a point forecast of curtailment is of limited use for decision-making. Curtailment is a zero-inflated, heavy-tailed quantity: on a large fraction of days a given plant curtails nothing, and on the days it does the amounts are highly variable. A single predicted number conveys neither the probability that curtailment will occur nor the range it might take. What a decision-maker needs is a calibrated statement of uncertainty, an interval that contains the realised value with a stated and reliable probability. Providing such intervals, and keeping them informative when the system changes regime, is the problem we address.

This paper makes two contributions. The first is a data contribution. We construct and release an open, plant-level dataset of curtailment in the SEN at daily and hourly resolution, spanning 53 continuous months and 300 plants across solar, wind and two hydropower technologies, built from the public reports of the system operator and reconciled against the official aggregate figures. In the course of building it we identified and corrected eight classes of errata in the original public files, using deterministic and reproducible rules, and recovered two calendar days that were missing from the operator's own daily breakdown. To our knowledge no open curtailment dataset at this granularity exists for a Latin American power system. The second is a methodological contribution. Using this dataset, we build a two-part probabilistic forecaster for plant-level curtailment and wrap it in a conformal calibration layer scored by the probability integral transform. We characterise how this regime change inflates standard calibration into over-covered, unnecessarily wide intervals, we quantify how much of the apparent change is seasonal and how much is secular, we show that the textbook remedy for covariate shift fails on a controlled synthetic panel by returning uninformative intervals, and we propose an adaptive scheme that restores interval sharpness through the transition while keeping coverage near nominal, distinguishing throughout between the finite-sample guarantee of the static split and the long-run guarantee of the adaptive layer.

The remainder of the paper is organised as follows. Section~\ref{sec:related} reviews related work on curtailment forecasting and on conformal prediction under distribution shift. Section~\ref{sec:data} describes the dataset and its construction. Section~\ref{sec:methods} presents the base forecaster and the conformal calibration layer. Section~\ref{sec:results} reports the empirical results on real data. Section~\ref{sec:discussion} discusses operational implications and threats to validity, and Section~\ref{sec:conclusion} concludes.

%==================================================================
\section{Related work}
\label{sec:related}

Curtailment forecasting has received far less attention than renewable generation forecasting, and existing studies have largely relied on point-forecast machine-learning models for individual grids, without addressing either distribution-free uncertainty quantification or structural regime change. The present work differs on both counts, and connects the applied problem to the literature on conformal prediction under distribution shift.

Conformal prediction constructs prediction sets with finite-sample, distribution-free coverage under exchangeability~\cite{angelopoulos2023}, the very assumption a regime change violates. Two lines of work relax it. Weighted conformal prediction extends the guarantee to covariate shift by reweighting calibration points according to a density ratio between the calibration and test distributions~\cite{tibshirani2019}; this presupposes that the conditional law $\mathbb{P}(Y\mid X)$ is invariant, which fails when the change is in the conditional itself. Conformal validity beyond exchangeability, with fixed non-uniform weights and an explicit coverage-gap bound in total variation, covers the regime-change case directly~\cite{barber2023}, and online adaptive conformal inference attains long-run coverage under arbitrary shifts by updating the miscoverage level from observed errors~\cite{gibbs2021}. Our adaptive scheme builds on these two, scores observations by their probability integral transform following the distributional-conformal construction of Chernozhukov et al.~\cite{chernozhukov2021}, and uses monotone one-dimensional optimal transport~\cite{peyre2019} to describe and correct the displacement of the score distribution across regimes. Proper scoring rules, in particular the continuous ranked probability score~\cite{gneiting2007}, provide the evaluation basis for the probabilistic forecasts.

%==================================================================
\section{Data}
\label{sec:data}

\subsection{Source and scope}

The dataset consolidates the curtailment records published monthly by Chile's National Electricity Coordinator (Coordinador El\'ectrico Nacional, CEN) into two validated relational tables. The daily table contains 293{,}678 records and the hourly table 6{,}906{,}761 records, both covering the continuous period from 1 January 2022 to 31 May 2026, that is 1{,}612 days built from all 53 monthly reports, with no missing days. Records span 300 generation plants across four technologies: solar photovoltaic (75 plants), wind (56), run-of-river hydropower (140) and reservoir hydropower (29). Twelve hydropower plants are labelled inconsistently across the operator's monthly reports, alternating between the run-of-river and reservoir categories; the plant registry resolves each of them to a single canonical technology, so that a count of distinct plant-technology pairs in the raw daily table returns 312 rather than 300. Hydropower curtailment records are present from June 2024 onward, with hourly hydropower detail from July 2024, following the coverage of the underlying public reports.

Each plant is linked to a metadata registry matched to the CEN installation catalogue, with 18 fields including the official identifier, installed capacity, owner, region, province, municipality, geographic coordinates, commissioning date, connection point and regulatory classification; the matching method is recorded for every plant. Days and hours without curtailment are stored explicitly as zero-valued records, which represent 49.8\% of daily and 83.6\% of hourly rows. This explicit encoding of zeros is deliberate: it preserves the zero-inflated structure of the quantity, which is essential for the modelling in Section~\ref{sec:methods}. Figure~\ref{fig:map} shows the geographic distribution of the plants, with the concentration of solar capacity in the northern regions and of hydropower in the south.

\begin{figure}[tp]
\centering
\includegraphics[width=0.45\textwidth]{fig1_plants_map.pdf}
\caption{Geographic distribution of the 300 power plants in the dataset, coloured and shaped by technology, with marker area proportional to each plant's cumulative curtailed energy over January 2022 to May 2026.}
\label{fig:map}
\end{figure}

\subsection{Construction and validation}

The 53 monthly reports were retrieved from the CEN public portal with a versioned download script, parsed and loaded into a relational database, and the original files were preserved unmodified with SHA-256 checksums. During construction we identified eight classes of errata in the original publications. Five were corrected or recovered deterministically in the curated tables and three are limitations of the source that we document without correction. Among the corrections, two omitted calendar days were recovered from the corresponding monthly hourly sheets and cross-checked to the millesimal against the source files' own total columns; one of these was detectable only by a structural continuity check, since the omitted day carried zero energy and was therefore invisible to any energy-based validation.

Two independent checks establish the fidelity of the result. Annual totals from the daily table match the CEN official summary exactly, to a difference of 0.000~GWh, in every year from 2022 to 2026, and aggregate 2022 wind and solar curtailment reproduces the 1.47~TWh independently reported by Fraunhofer Chile for that year. At monthly resolution the agreement is close but not exact: in a census month-by-month reconciliation of hourly sums against daily totals over all 53 months, 50 months agree within 1\% and 35 agree exactly, while the three remaining months correspond to documented inconsistencies within the operator's own publications. The residual sub-percent differences arise from the operator reassigning the same monthly energy differently across sister run-of-river units between its monthly hourly sheets and its annual closing, with a net effect of zero at system level, which is the same mechanism as errata classes six and seven. The dataset totals 19{,}030~GWh of recorded curtailment over the period.

The dataset is publicly available on Zenodo under a CC~BY~4.0 licence~\cite{dataset2026}, together with the data dictionary, the errata log and the integrity checksums. The processing pipeline that regenerates every curated table is available in the companion code repository.

\subsection{Exploratory characterisation}

Three features of the data govern the modelling choices that follow. First, curtailment is strongly zero-inflated and the occurrence of curtailment is persistent in time, with a daily lag-one autocorrelation of roughly 0.85 at system level and a marked weekly cycle. Second, the positive magnitudes are heavy-tailed and approximately lognormal, with a ratio between the 99th and 50th percentiles of the order of twenty-five across technologies, and of about fifteen for solar alone. Third, the change in the distribution of curtailment over the study period acts on scale, not on occurrence and not on timing. The occurrence probability is stable: for solar plants in the northern regions it stays within 0.71 and 0.85 throughout 2024 to 2026, and a change-point search on the monthly series of $\mathbb{P}(Y>0)$ returns no break after September 2022. The intraday shape is equally stable: the peak hour of the normalised intraday profile is hour 16 in every half-year of the record, the centroid of the profile moves by less than fifteen minutes between half-years, and the total variation distance between the profiles of consecutive half-years stays between 0.04 and 0.06, never above the 0.08 measured between 2023-H2 and 2024-H1, a pair across which no regime change is postulated. The positive magnitudes, in contrast, roughly triple: the median of positive daily curtailment for solar plants in Antofagasta and Atacama rises from 92~MWh in the calibration window to 267~MWh in October to December 2024. Figure~\ref{fig:dist} illustrates the first two features. These three facts, zero inflation, heavy tails and a scale-type shift, motivate respectively the two-part model, the choice of evaluation metric and the adaptive calibration described next.

\begin{figure}[tp]
\centering
\includegraphics[width=0.75\textwidth]{fig2_distribution.pdf}
\caption{Distributional structure of plant-level curtailment. Top: share of zero-valued plant-days by technology. Bottom: distribution of positive magnitudes on a logarithmic scale, with the 50th, 90th and 99th percentiles marked. The combination of a large point mass at zero and a heavy right tail motivates the two-part model of Section~\ref{sec:methods}.}
\label{fig:dist}
\end{figure}

The timing of that change in scale is best read after removing the annual cycle. Figure~\ref{fig:regime} shows the monthly mean of $\log Y$ over positive plant-days for solar plants in Antofagasta and Atacama, together with the same series net of a calendar-month median. Applied to the deseasonalised series under an identical BIC-type penalty, the PELT and binary-segmentation algorithms return two breaks and agree on both: May 2023, with a jump of $+1.05$ in mean log-magnitude, and January 2024, with a jump of $+0.76$. Neither algorithm returns a break in October 2024, and neither returns any break in 2025. The change is therefore a ramp of roughly eighteen months rather than a step, and it is essentially complete by the end of 2024: comparing like with like across years, the Wasserstein-1 distance between the positive log-magnitudes of October to December 2023 and those of October to December 2024 is 1.20, whereas between October to December 2024 and October to December 2025 it is 0.085, below the 0.105 threshold of a permutation null that resamples whole days and therefore not distinguishable from no change at all.

\begin{figure*}[tp]
\centering
\includegraphics[width=\textwidth]{fig5_regime_changepoints.pdf}
\caption{Monthly mean of $\log Y$ over positive plant-days, solar plants in Antofagasta and Atacama. Top: raw series, dominated by the austral annual cycle. Bottom: the same series net of a calendar-month median, with the two change points returned by PELT and by binary segmentation marked by vertical lines (May 2023 and January 2024) and the calibration and transition windows of Section~\ref{sec:methods} shaded. Neither algorithm returns a break in October 2024 or in 2025; the deseasonalised level reaches its maximum in October and November 2024 and then stops rising. The isolated spike of June 2025, discussed in Section~\ref{sec:discussion}, is the largest deseasonalised anomaly of the post-calibration record.}
\label{fig:regime}
\end{figure*}

Because the annual cycle is strong, any contrast between windows that fall in different seasons registers a large distance even when no regime change is present, and this bears directly on how the evaluation windows of Section~\ref{sec:methods} should be read. Table~\ref{tab:seasonal} decomposes the January--September to October--December shift in the monthly mean of $\log Y$ into a seasonal component, given by the calendar-month effects and therefore identical in every year, and a secular remainder. In 2024 the shift is the largest of the record, but 61\% of it for northern solar, and 76\% at system level, is the seasonal component that recurs every year; the same contrast yields a Wasserstein-1 distance of 0.96 in 2023 and 0.88 in 2025, years for which no regime change is claimed. The secular remainder of 2024, $+0.46$ for northern solar, is real and is the largest of the four years, but it is of the same order as the $+0.29$ of 2023. We report the raw contrast throughout the results, and flag here that it should not be read as if it were entirely a regime effect.

\begin{table}[t]
\centering
\caption{Decomposition of the January--September to October--December shift in the monthly mean of $\log Y$ over positive plant-days, in log units, by year. The seasonal component is the difference between the calendar-month effects of the two windows and is by construction the same in every year; the secular component is the remainder. The 2024 row is the contrast that defines the transition window used in the evaluation.}
\label{tab:seasonal}
\begin{tabular}{lccc}
\toprule
Year & Total shift & Seasonal & Secular \\
\midrule
\multicolumn{4}{l}{\emph{Solar, Antofagasta and Atacama}} \\
2022 & $+0.64$ & $+0.72$ & $-0.09$ \\
2023 & $+1.01$ & $+0.72$ & $+0.29$ \\
\textbf{2024} & $\mathbf{+1.18}$ & $\mathbf{+0.72}$ & $\mathbf{+0.46}$ \\
2025 & $+0.90$ & $+0.72$ & $+0.17$ \\
\midrule
\multicolumn{4}{l}{\emph{System (solar and wind)}} \\
2022 & $+0.77$ & $+0.66$ & $+0.11$ \\
2023 & $+0.89$ & $+0.66$ & $+0.23$ \\
\textbf{2024} & $\mathbf{+0.86}$ & $\mathbf{+0.66}$ & $\mathbf{+0.20}$ \\
2025 & $+0.71$ & $+0.66$ & $+0.05$ \\
\bottomrule
\end{tabular}

\vspace{3pt}
\parbox{\linewidth}{\footnotesize Note: the calendar-month effects are estimated as the median of each calendar month over the whole record. On a 53-month series carrying a strong trend, an STL-type decomposition would return somewhat different values; it would not change the order of magnitude of the split, and the secular components would remain the smaller part of the 2024 contrast.}
\end{table}

%==================================================================
\section{Methods}
\label{sec:methods}

\subsection{Problem setup and base forecaster}

For each plant and each origin day $t$ we predict curtailment at horizon $t+7$. Because roughly half of plant-days carry no curtailment and the positive amounts are heavy-tailed, a single regression is ill-suited and we adopt a two-part (hurdle) model. Let $x$ denote the feature vector available at the origin, built only from information observed up to $t$, comprising lagged curtailment, moving averages, a lagged occurrence indicator, calendar terms and static plant attributes. The first part is a classifier estimating the occurrence probability $p(x)=\mathbb{P}(Y>0\mid x)$. The second part models the positive magnitude as lognormal, $\log Y \mid Y>0, x \sim \mathcal{N}(\mu(x),\sigma^2)$. The predictive law is the mixture
\begin{equation}
Y \mid x \;\sim\; \big(1-p(x)\big)\,\delta_0 \;+\; p(x)\,\mathrm{LogNormal}\big(\mu(x),\sigma^2\big),
\label{eq:hurdle}
\end{equation}
where $\delta_0$ is a point mass at zero. Writing $\Phi$ for the standard normal CDF, the mixture CDF is $F_x(y)=0$ for $y<0$, has an atom $F_x(0)=1-p(x)$ at the origin (so that $F_x(0^-)=0$), and for $y>0$ equals
\begin{equation*}
F_x(y)=\big(1-p(x)\big)+p(x)\,\Phi\!\Big(\tfrac{\log y-\mu(x)}{\sigma}\Big).
\end{equation*}
Both parts are fitted with gradient boosting on data up to the end of 2023. The dispersion $\sigma$ is estimated out of fold to avoid the downward bias of the in-sample residual variance of a boosted model; on the real data $\hat\sigma=1.70$, which we return to in Section~\ref{sec:results} because it sets the sharpness floor of the whole pipeline.

\subsection{Conformal calibration with a PIT score}

A point forecast, or even the raw mixture~\eqref{eq:hurdle}, does not by itself guarantee that a stated interval contains the realised value with the nominal probability. We obtain that guarantee by conformal calibration. Rather than measuring nonconformity in energy units, which are not comparable across plants of very different size, we score each observation by its probability integral transform (PIT) under the predictive law,
\begin{equation}
s \;=\; F_x(y),
\label{eq:pit}
\end{equation}
following the distributional-conformal construction of Chernozhukov et al.~\cite{chernozhukov2021}. Because $F_x$ has an atom at zero, the score is randomised there in the standard way for mixed distributions~\cite{smith1985}: for $y=0$ we set $s=F_x(0^-)+U\big(F_x(0)-F_x(0^-)\big)=\big(1-p(x)\big)\,U$ with $U\sim\mathrm{Uniform}(0,1)$, while for $y>0$ no randomisation is needed. Under the true predictive law $s$ is then $\mathrm{Uniform}(0,1)$, and it lives on a common $[0,1]$ scale across plants of any size, which is why it handles the atom and the plant heterogeneity cleanly.

This score has a further property that is decisive under regime change. For a positive observation the score is a monotone transform of the standardised log-magnitude $z=(\log y-\mu(x))/\sigma$, namely $s=(1-p)+p\,\Phi(z)$. A change of scale of the magnitude distribution by a factor $c$ is a location shift $\log c$ of $\log Y$, hence a translation of $z$ by $\delta=\log c/\sigma$. It is therefore a \emph{translation on the latent Gaussian (probit) scale}, and only a monotone displacement of the law of $s$ on $[0,1]$; it is not a translation of $s$ itself. This low-dimensional structure, established under the empirically supported assumption that the shift rescales the positive magnitudes while leaving the occurrence probability and the intraday timing essentially unchanged (Section~\ref{sec:data}), is what lets the displacement be estimated from a short recent window while the shape of the tail is inherited from the large calibration set.

Given a calibration set of size $n$ scored by~\eqref{eq:pit}, the conformal quantile is the order statistic $\hat q=s_{(k)}$ with $k=\lceil (n+1)(1-\alpha)\rceil$, which yields the finite-sample marginal guarantee $\mathbb{P}(Y \le U(x)) \ge 1-\alpha$ under exchangeability; when $k>n$ the quantile is $+\infty$ and the interval is the whole line. The upper prediction limit $U(x)=F_x^{-1}(\hat q)$ is recovered by inverting the mixture: it is zero when $\hat q\le 1-p(x)$ (the quantile falls in the atom), it is $+\infty$ when $\hat q\ge 1$, and otherwise
\begin{equation}
U(x) \;=\; \exp\!\Big(\mu(x) + \sigma\,\Phi^{-1}\!\big(\textstyle\frac{\hat q-(1-p(x))}{p(x)}\big)\Big).
\label{eq:invert}
\end{equation}
The two boundary cases are the two ends of the clipped argument $\big(\hat q-(1-p)\big)/p\in[0,1]$ in~\eqref{eq:invert}. A short proof of the inversion is given in \ref{app:inversion}, and it was additionally checked numerically against oracle coverage.

\subsection{Coverage under regime change}

Exchangeability fails when the system changes regime, and the distributional shift documented in Section~\ref{sec:data} is such a failure. It fails, moreover, already within each regime: the daily systemic state is shared by all plants, so the calibration scores arrive in strongly correlated daily blocks, and the effective sample size for the fluctuation of coverage is of the order of the number of days, not of rows. For this reason every coverage figure we report carries a standard error clustered by date. We consider four responses to the shift.

\emph{Weighted conformal.} Reweighting each calibration point by an estimate of the density ratio between calibration and test distributions is the standard remedy for \emph{covariate} shift~\cite{tibshirani2019}. It is the wrong tool here: the shift changes the conditional law $\mathbb{P}(Y\mid x)$, not the marginal of $x$, so no reweighting of old scores can reflect the new conditional, and the density-ratio estimate degenerates because features that are monotone in time separate calibration from test almost perfectly. On the synthetic panel, where covariate shift can be isolated, the per-point implementation of~\cite{tibshirani2019} returns infinite intervals for a large fraction of test points (Section~\ref{sec:results}); we report this as a negative result that motivates the alternatives, and we take the fixed-recency-weight scheme of Barber et al.~\cite{barber2023}, which bounds the coverage gap in total variation without a classifier, as the correct recency benchmark.

\emph{Sliding-window recalibration.} We discard old calibration data and recalibrate on a trailing 60-day window refreshed every 7 days. Because the base model forecasts at a 7-day horizon, the label of an origin is realised only 7 days later; to avoid using information unavailable at the forecast origin, every online calibration window ends 7 days before its target (a horizon embargo). The static split does not use online calibration and is therefore unaffected by the embargo, which serves as an implementation control.

\emph{Adaptive conformal inference (ACI).} We update the effective miscoverage level online~\cite{gibbs2021}, $\alpha_{t+1}=\alpha_t+\gamma(\alpha-\mathrm{err}_t)$ with $\mathrm{err}_t=\mathbb{1}\{Y_t>U_t\}$, where the error feedback is delayed by the same 7-day embargo. This yields long-run average coverage without exchangeability, at the cost of the finite-sample guarantee of the static split.

\emph{Transport of scores.} Motivated by the monotone-displacement structure above, we transport the calibration scores from the old regime onto the new one before conformalising. Let $\hat F_{\mathrm{old}}$ and $\hat F_{\mathrm{new}}$ be the empirical CDFs of the calibration scores and of the scores in the recent window; the monotone one-dimensional optimal-transport map is the increasing rearrangement $T=\hat F_{\mathrm{new}}^{-1}\circ\hat F_{\mathrm{old}}$, estimated nonparametrically by interpolation between empirical quantiles. In the upper tail, where the short window is noisy, the map is shrunk towards the identity in proportion to a bootstrap estimate of its variance, so that a noisy new quantile is replaced by the stable old one rather than extrapolated. Because the transport map is data-dependent it breaks the exchangeability of the transported scores with the test score; we therefore restore validity by wrapping the transported scores in the online adaptive scheme, so that coverage is controlled by the ACI update and the transport contributes only sharpness. An idealised validity argument, exact under an oracle map and controlled through the total-variation framework of~\cite{barber2023} under an estimated map, is given in \ref{app:coverage}.

\subsection{Evaluation}

We evaluate probabilistic forecasts by empirical coverage against the nominal level, by mean interval width over finite intervals (sharpness), by the fraction of infinite intervals, and by the continuous ranked probability score (CRPS), which is a property of the base predictive rather than of the calibration layer and is therefore reported once per period. Because all plants share the state of the system on a given day, the standard errors of every coverage figure are clustered by date; these are of the order of one percentage point, so that differences between methods smaller than a few points are within noise. The calibration set is January to August 2024. The first test month is September 2024. The following window, October to December 2024, is the transition window on which the calibration layer is stressed, and it is selected as the window of maximum divergence from the calibration distribution rather than as the interval of entry of any particular technology: measured by the Wasserstein-1 distance between the positive log-magnitudes of each test window and those of the calibration set, it scores 1.25 against 0.28 to 0.76 for every other window. That distance is almost entirely a displacement of location: the change is close to a pure translation of the distribution of log-magnitudes, and for a pure translation the Wasserstein-1 distance equals the difference in means, the two agreeing to within 0.001 for this contrast. Table~\ref{tab:seasonal} decomposes the analogous January--September to October--December contrast, which is defined on a like-for-like basis so that the same four years can be compared; carried out instead on the calibration window itself, January to August, the split is essentially the same, with 61\% of the displacement seasonal for northern solar and 74\% at system level. On that basis the window is best understood as the hardest case the calibration layer faces in this record, not as the signature of a single event. The remaining periods are taken by semester. The base forecaster is benchmarked against persistence, a seasonal-naive predictor, a point gradient-boosting model and a quantile gradient-boosting model.

%==================================================================
\section{Results}
\label{sec:results}

\subsection{Point forecasting offers little marginal value}

Table~\ref{tab:mae} reports the mean absolute error of the base models on the out-of-sample period. The seasonal-naive predictor, which simply averages the same weekday over recent weeks, is the most accurate of all, and is not beaten by the gradient-boosting models; the hurdle mixture median, which is optimised for the full predictive law rather than for absolute error, is the least accurate on this metric. This ordering is a direct consequence of the strong weekly persistence documented in Section~\ref{sec:data}: when a quantity is that persistent, sophisticated point models add little over a well-chosen calendar rule. The operational conclusion is that value does not lie in improving the predicted number, but in quantifying its uncertainty reliably, which is the object of the conformal layer.

\begin{table}[t]
\centering
\caption{Mean absolute error (MWh) of the base forecasters on the out-of-sample period. Lower is better.}
\label{tab:mae}
\begin{tabular}{lcccc}
\toprule
Model & 2024 & 2025 & 2026 & Total \\
\midrule
Seasonal-naive (weekly)      & 84.3 & 86.3 & 84.2 & 85.1 \\
Quantile GBM (median)        & 88.5 & 89.9 & 84.3 & 88.3 \\
Persistence                  & 88.4 & 96.3 & 90.7 & 92.2 \\
Point XGBoost                & 92.7 & 94.7 & 89.7 & 93.0 \\
Hurdle (mixture median)      & 104.4 & 104.8 & 101.6 & 104.1 \\
\bottomrule
\end{tabular}
\end{table}

\subsection{Calibration under the regime change}

On real data the shift does not break coverage; it inflates it. Under a static conformal split, empirical coverage stays between 89\% and 95\% in every period (Table~\ref{tab:results}), so the severe over-coverage seen on synthetic data is absent, because the base forecaster, through its lagged and moving-average features, absorbs much of the shift into its own predictive. The cost of the shift is therefore paid in width rather than in coverage, and it is concentrated in the transition window: the static split reaches 95.3\% coverage with a mean width of 1203~MWh during October to December 2024. Part of that widening is seasonal rather than secular, in the proportion given in Table~\ref{tab:seasonal}, and the same qualitative pattern of widening would be expected in the fourth quarter of any year in this record; what distinguishes 2024 is that the seasonal peak lands on top of the largest secular displacement of the four years. Figure~\ref{fig:rolling} traces the behaviour over time: the static split drifts above the nominal level as the window advances, while the adaptive schemes track it more closely.

\begin{figure*}[t]
\centering
\includegraphics[width=\textwidth]{fig3_rolling_coverage.pdf}
\caption{Rolling empirical coverage (90-day window) of the conformal methods on real data, computed with the 7-day horizon embargo. The dashed line marks the 90\% nominal level and the shaded band marks the window of maximum divergence from the calibration distribution, October to December 2024. The static split drifts into over-coverage as that window advances, whereas the adaptive schemes remain close to nominal.}
\label{fig:rolling}
\end{figure*}

The adaptive schemes pay off precisely there. Transport of the calibration scores wrapped in online adaptation ($\gamma=0.05$) brings the transition intervals down to 806~MWh at 90.9\% coverage, a reduction of about a third at a coverage within one point of nominal, and the sliding window gives 834~MWh at 92.0\%. Outside that window the trade-off is ambiguous: with the static split already near nominal, the adaptive methods mostly exchange coverage for width and can under-cover, as in 2025-S1, and their gains are small and within the date-clustered standard error. This is what the data predict rather than a shortcoming of the adaptive layer. The first semester of 2025 is not distinguishable from the pre-transition regime: the Wasserstein-1 distance between the positive log-magnitudes of January to September 2024 and those of January to June 2025 is 0.171, with a 95\% bootstrap interval of $[0.068, 0.377]$ resampling whole days, against a threshold of 0.215 under a permutation null that also resamples whole days. There is no shift left to correct in that period, so an adaptive layer can only trade coverage for width, and that it does so by margins within the date-clustered standard error is evidence that it does not damage calibration when adaptation is not called for. Figure~\ref{fig:tradeoff} makes this trade-off explicit across methods and periods. Two honest limitations frame the result. First, the base dispersion is large on real data ($\hat\sigma=1.70$), so a substantial part of every width originates in the base model rather than in the calibration layer, which calibrates that dispersion rather than removing it; the sharpness floor is a property of the hurdle, not of the conformal step. Second, the weighted-conformal failure that motivates the adaptive schemes is established on the synthetic panel, where covariate shift can be isolated: the per-point implementation returns infinite intervals for the majority of test points during the transition. It is a negative result about the tool, not about the data, and it is reported as such.

\begin{table*}[t]
\centering
\footnotesize
\caption{Coverage and mean interval width of the conformal methods by period, on real data, with the 7-day horizon embargo. Each cell shows coverage in percent (clustered standard error by date in parentheses) over mean interval width in MWh. Nominal coverage is 90\%. The base-model CRPS is a property of the predictive, not of the calibration layer, and is reported once per period. Superscripts give the fraction of infinite intervals where nonzero.}
\label{tab:results}
\begin{tabular}{lccccc}
\toprule
Method & test\_pre & test\_ramp & 2025-S1 & 2025-S2 & 2026-S1 \\
       & (Sep 24, stable) & (Oct--Dec 24, max.\ divergence) & & & \\
\midrule
Static split (PIT)      & 93.5\,(1.03) / 799 & 95.3\,(0.45) / 1203 & 89.4\,(1.04) / 455 & 94.0\,(0.45) / 750 & 92.5\,(0.58) / 535 \\
Sliding window 60d      & 93.0\,(1.04) / 742 & 92.0\,(0.73) / 834  & 85.4\,(1.16) / 323 & 92.9\,(0.51) / 650 & 89.5\,(0.76) / 419 \\
ACI ($\gamma{=}0.02$)   & 93.0\,(1.06) / 759 & 93.5\,(0.55) / 915  & 85.7\,(1.16) / 345 & 93.5\,(0.49) / 669 & 89.1\,(0.76) / 405 \\
ACI ($\gamma{=}0.05$)   & 92.4\,(1.16) / 718 & 92.1\,(0.67) / 822  & 86.6\,(1.14) / 405\textsuperscript{2.8} & 92.5\,(0.58) / 609\textsuperscript{11.8} & 89.8\,(0.75) / 446 \\
Transport+ACI ($\gamma{=}0.02$) & 92.6\,(1.07) / 713 & 91.1\,(0.80) / 787 & 86.5\,(1.17) / 437\textsuperscript{0.6} & 92.4\,(0.60) / 986\textsuperscript{1.1} & 90.3\,(0.73) / 459 \\
Transport+ACI ($\gamma{=}0.05$) & 92.0\,(1.19) / 681 & 90.9\,(0.81) / 806 & 87.4\,(1.13) / 446\textsuperscript{9.9} & 91.6\,(0.63) / 718\textsuperscript{10.7} & 90.5\,(0.72) / 489 \\
\midrule
Base-model CRPS (MWh)   & 96.1 & 133.5 & 70.5 & 94.5 & 79.0 \\
\bottomrule
\end{tabular}

\vspace{3pt}
\parbox{\textwidth}{\footnotesize Note: the key \texttt{test\_ramp} is retained from the code and from the published result files for reproducibility; the window itself is defined as the one of maximum divergence from the calibration distribution, not by any storage deployment.}
\end{table*}

\begin{figure*}[t]
\centering
\includegraphics[width=\textwidth]{fig4_coverage_width.pdf}
\caption{Coverage and sharpness by method and period. Left: empirical coverage with date-clustered standard errors against the 90\% nominal level. Right: mean interval width. Asterisks mark bars whose mean width is computed over finite intervals only, because the corresponding configuration returns at least 1\% of infinite intervals. The adaptive schemes reduce width during the transition window while holding coverage near nominal.}
\label{fig:tradeoff}
\end{figure*}

%==================================================================
\section{Discussion}
\label{sec:discussion}

The results carry a clear operational message. For plant-level curtailment at a seven-day horizon, effort is better spent on reliable uncertainty quantification than on chasing point accuracy, which a calendar baseline already delivers. The value of the method is not that it predicts curtailment better in absolute terms, but that it attaches to each prediction an interval that remains informative even as the system changes regime. This is what a decision-maker who manages a portfolio of assets actually needs.

Three threats to validity deserve emphasis. First, exchangeability is violated even within a single regime, because the daily systemic state is shared across plants and the scores arrive in correlated daily blocks; the finite-sample guarantee of the static split is therefore marginal and approximate, and this is why every coverage figure is reported with a date-clustered standard error. Second, coverage differences of one to three points between methods lie within that standard error, and we refrain from over-interpreting small gaps. Third, because a decision-maker evaluates each plant individually, conditional reliability matters: a method that is sharper on average but less reliable on the largest plants, which carry the greatest economic weight, is not necessarily preferable. Fourth, and conceptually, the record contains two different kinds of change and the calibration layer responds to only one of them. A displacement of the level of the predictive law, such as the one that accumulates through 2023 and 2024, is what the transport map and the online update are built to absorb, because it appears as a displacement of the score distribution. A break in the rate of growth is not. Between December 2024 and January 2025 the year-over-year growth of the monthly mean log-magnitude falls from $+0.89$ to $+0.09$ per year for northern solar and from $+0.90$ to $+0.01$ at system level, while the northern solar fleet keeps expanding by 8 per cent through 2025 and curtailment per installed MW in those regions falls by 15 per cent in the fourth quarter. None of the methods evaluated here detects that break, and none needs to in order to stay calibrated, because a trend that flattens leaves no displacement to correct. It matters for interpretation nonetheless: a calibration layer that is reliable under shifts of level carries no guarantee under changes of trend, and the two are not distinguishable from coverage alone.

The coverage dip of mid-2025 can now be given a date and a signature. It is concentrated in June 2025, the largest deseasonalised monthly anomaly of the entire post-calibration record: the monthly mean log-magnitude for northern solar stands $+1.06$ above the deseasonalised level of the calibration window, against a maximum of $+0.67$ for any single month of the October to December 2024 window, and the Wasserstein-1 distance between May to July 2025 and August to December 2025 is 1.335, larger than the 1.146 of the transition contrast itself. Its signature is the reverse of the 2024 window: magnitudes rise while the occurrence probability falls, from 0.755 in June 2024 to 0.623 in June 2025 at system level, whereas in October to December 2024 the two rose together. It appears in solar and in wind separately and across Antofagasta, Atacama, Coquimbo and Biob\'io, which rules out a single plant or a residual errata, and the operator's monthly report for that month records no storage system entering operation. A one-month excursion with an occurrence signature opposite to that of the main transition is the kind of event that a continuously adapting scheme absorbs only slowly and that a change-point trigger would catch, which strengthens the case for the detector proposed below.

On the causes of the change we are deliberately agnostic, and the dataset released here cannot settle the question, because it records curtailment and not the assets that absorb it. The entry of battery storage is one candidate for the flattening of the trend at the turn of 2024 to 2025 and is quantitatively of the right size, but the evidence is bounded on three sides at once: comparing October to December 2025 with the same quarter of 2024 and holding the solar fleet constant, 257~GWh of curtailment did not occur, or 2.80~GWh per day, against a growth of 3.10~GWh in operating storage energy in Antofagasta and Atacama over the same interval, which covers the gap only under the upper-bound assumptions of one full daily cycle, unit efficiency and every charged megawatt-hour having otherwise been curtailed; the same exercise repeated for January to May 2026 fails to reproduce it, since storage energy grew by a further 3.43~GWh while the gap was only 18~GWh; and at least three alternative mechanisms operate over the same months and cannot be separated from storage with these data, namely reinforcements of the north-to-south transmission corridor, changes in northern mining demand, and the mechanical saturation of curtailment itself once the corridor is full throughout the solar window. We therefore record storage as one bounded hypothesis among several and make no causal claim. A registry of storage systems with their commercial operation dates, provided in the supplementary material, is a prerequisite for testing it properly.

Several directions remain open. A change-point detector could trigger recalibration at the onset of a shift rather than adapting continuously, which would address the June 2025 excursion described above; the same machinery used offline in Section~\ref{sec:data} to date the breaks is a natural starting point for an online version. A heteroscedastic base model with a conditional dispersion $\sigma(x)$ was tested and yielded sharper intervals on average but degraded conditional coverage on the largest plants; the diagnosis is that the shift acts on the location rather than on the dispersion of the predictive, so modelling the variance cannot correct a shift in the mean, and we retain the constant-dispersion model for operational reliability. Extending the transport-based calibration to the multivariate setting, across plants that share a transmission corridor, is a natural next step.

%==================================================================
\section{Conclusion}
\label{sec:conclusion}

We have presented an open, errata-corrected, plant-level dataset of renewable curtailment for Chile's National Electricity System and a probabilistic forecasting method that pairs a two-part model with a conformal calibration layer designed to remain informative under regime change. The empirical study shows that point forecasting offers little marginal value for this problem, that a regime shift of the kind observed here inflates standard conformal intervals into over-covered and wider bands rather than breaking their coverage, that a substantial share of what reads as regime change is the annual cycle and has to be separated from it before the shift can be sized, that the textbook remedy for covariate shift fails on a controlled panel by returning uninformative intervals, and that an adaptive scheme based on a probability-integral-transform score with online recalibration and score transport restores sharpness through the transition while holding coverage near nominal. The dataset and code are publicly available, and we hope both the resource and the method are useful to operators, asset owners and researchers working on the integration of variable renewable energy.

%==================================================================
\section*{CRediT authorship contribution statement}
\textbf{Pablo Reyes Cerda:} Conceptualization, Data curation, Software, Validation, Investigation, Writing -- original draft. \textbf{Kerven Cea Morales:} Methodology, Formal analysis, Validation, Writing -- review and editing.

\section*{Declaration of competing interest}
Pablo Reyes Cerda is the founder of Synapta SpA and of CurtailmentIQ, a commercial forecasting service that uses data sources related to those described in this article. Kerven Cea Morales declares no competing interest.

\section*{Declaration of generative AI in scientific writing}
During the preparation of this work the authors used an AI-based assistant to improve the readability and language of the manuscript. After using this tool, the authors reviewed and edited the content as needed and take full responsibility for the content of the article.

\section*{Data availability}
The dataset is openly available on Zenodo under a CC~BY~4.0 licence (DOI: 10.5281/zenodo.21198816, which resolves to the latest version of the deposit). The processing pipeline, the modelling code and the scripts that regenerate every table and figure in this article are available in the companion repository at \url{https://github.com/Pablosebastian-reyes/curtailmentiq-research}. The supplementary material accompanying this article contains a registry of the energy storage systems of the SEN with their commercial operation dates and sources, compiled from the public reports of the system operator, the Ministry of Energy and the National Energy Commission; it is offered so that the bounded hypothesis of Section~\ref{sec:discussion} can be tested independently, and it is not part of the released dataset.

%==================================================================
\appendix

\section{Inversion of the PIT limit}
\label{app:inversion}

Fix $x$ and write $p=p(x)$, $\mu=\mu(x)$. The mixture CDF is $F_x(y)=0$ for $y<0$, $F_x(0)=1-p$, and $F_x(y)=(1-p)+p\,\Phi\big((\log y-\mu)/\sigma\big)$ for $y>0$, which is continuous and strictly increasing on $(0,\infty)$ from $1-p$ to $1$. Given a level $\hat q\in[0,1]$ we seek the smallest $u\ge 0$ with $F_x(u)\ge \hat q$. If $\hat q\le 1-p$ the atom already reaches the level, so $u=0$. If $\hat q\ge 1$ no finite $u$ attains it and $u=+\infty$. For $1-p<\hat q<1$, strict monotonicity gives a unique solution of $(1-p)+p\,\Phi\big((\log u-\mu)/\sigma\big)=\hat q$, that is $\Phi\big((\log u-\mu)/\sigma\big)=(\hat q-(1-p))/p\in(0,1)$, hence $\log u=\mu+\sigma\,\Phi^{-1}\big((\hat q-(1-p))/p\big)$ and
\[
U(x)=\exp\!\Big(\mu+\sigma\,\Phi^{-1}\big(\tfrac{\hat q-(1-p)}{p}\big)\Big),
\]
which is~\eqref{eq:invert}. The two boundary cases are recovered as the limits of the clipped argument $(\hat q-(1-p))/p\to 0^+$ and $\to 1^-$, since $\Phi^{-1}(0^+)=-\infty$ gives $U=0$ and $\Phi^{-1}(1^-)=+\infty$ gives $U=+\infty$. \qed

\section{Coverage under score transport}
\label{app:coverage}

\emph{Static split.} If the calibration scores $s_1,\dots,s_n$ and the test score $s_{n+1}$ are exchangeable, then with $\hat q=s_{(k)}$, $k=\lceil(n+1)(1-\alpha)\rceil$, one has $\mathbb{P}(s_{n+1}\le\hat q)\ge k/(n+1)\ge 1-\alpha$, and since $F_x$ is monotone, $s_{n+1}\le\hat q$ if and only if $Y\le U(x)$; hence $\mathbb{P}(Y\le U(x))\ge 1-\alpha$. This is a \emph{marginal} statement, and it is exact only under exchangeability; the serial and cross-sectional dependence of the panel make it approximate, which is why the reported standard errors are clustered by date.

\emph{Transported scores.} Let $P_{\mathrm{old}}$ and $P_{\mathrm{new}}$ be the laws of the score in the old and new regimes and $T=F_{\mathrm{new}}^{-1}\circ F_{\mathrm{old}}$ the monotone transport map. Under an \emph{oracle} map, $T(s)\sim P_{\mathrm{new}}$ for $s\sim P_{\mathrm{old}}$, so the transported calibration scores are exchangeable with the test score and the static guarantee transfers exactly. With an \emph{estimated} map $\hat T$ the transported scores are only approximately distributed as $P_{\mathrm{new}}$, and the coverage gap is bounded by the total-variation distance $d_{\mathrm{TV}}\big(\hat T_\# P_{\mathrm{old}},P_{\mathrm{new}}\big)$ through the weighted-exchangeability framework of Barber et al.~\cite{barber2023}. We do not rely on this bound in practice: because $\hat T$ is data-dependent and the gap is not directly controllable, we obtain validity operationally by feeding the transported scores into the adaptive update of Gibbs and Cand\`es~\cite{gibbs2021}, whose long-run coverage $\frac1T\sum_t\mathrm{err}_t\to\alpha$ holds for any sequence, so that the transport contributes sharpness and the online control contributes coverage.

%==================================================================
\begin{thebibliography}{99}

\bibitem{oshaughnessy2020}
E.~O'Shaughnessy, J.~R.~Cruce, K.~Xu,
Too much of a good thing? Global trends in the curtailment of solar PV,
Solar Energy 208 (2020) 1068--1077.

\bibitem{angelopoulos2023}
A.~N.~Angelopoulos, S.~Bates,
Conformal prediction: a gentle introduction,
Foundations and Trends in Machine Learning 16 (4) (2023) 494--591.

\bibitem{tibshirani2019}
R.~J.~Tibshirani, R.~Foygel~Barber, E.~Cand\`es, A.~Ramdas,
Conformal prediction under covariate shift,
Advances in Neural Information Processing Systems 32 (2019).

\bibitem{barber2023}
R.~Foygel~Barber, E.~J.~Cand\`es, A.~Ramdas, R.~J.~Tibshirani,
Conformal prediction beyond exchangeability,
The Annals of Statistics 51 (2) (2023) 816--845.

\bibitem{gibbs2021}
I.~Gibbs, E.~Cand\`es,
Adaptive conformal inference under distribution shift,
Advances in Neural Information Processing Systems 34 (2021).

\bibitem{gneiting2007}
T.~Gneiting, A.~E.~Raftery,
Strictly proper scoring rules, prediction, and estimation,
Journal of the American Statistical Association 102 (477) (2007) 359--378.

\bibitem{peyre2019}
G.~Peyr\'e, M.~Cuturi,
Computational optimal transport,
Foundations and Trends in Machine Learning 11 (5--6) (2019) 355--607.

\bibitem{chernozhukov2021}
V.~Chernozhukov, K.~W\"uthrich, Y.~Zhu,
Distributional conformal prediction,
Proceedings of the National Academy of Sciences 118 (48) (2021).
doi:10.1073/pnas.2107794118.

\bibitem{smith1985}
J.~Q.~Smith,
Diagnostic checks of non-standard time series models,
Journal of Forecasting 4 (3) (1985) 283--291.
doi:10.1002/for.3980040305.

\bibitem{dataset2026}
P.~Reyes~Cerda, K.~Cea~Morales,
An errata-corrected, plant-level dataset of solar, wind and hydropower curtailment in Chile's National Electricity System (2022--2026) [Data set],
Zenodo (2026). doi:10.5281/zenodo.21198816.

\end{thebibliography}

\end{document}
